\documentclass[10pt,letterpaper]{article}
\usepackage[margin=0.83in]{geometry}
\usepackage{mathptmx}
\usepackage[T1]{fontenc}
\usepackage{microtype}
\usepackage{enumitem}
\usepackage{titlesec}
\pagestyle{empty}
\setlength{\parindent}{1em}
\setlength{\parskip}{3pt}
\setlist[itemize]{leftmargin=1.25em,topsep=2pt,itemsep=3pt,parsep=0pt}
\titlespacing*{\section}{0pt}{5pt}{2pt}
\titleformat{\section}{\normalsize\bfseries}{}{0pt}{}
\begin{document}
\begin{center}
{\large\bfseries Knowledge-Grounded Reasoning and Self-Improvement for Autonomous AI Agents\par}
\vspace{2pt}
{\small Xinya Du, University of Texas at Dallas \quad | \quad ONR FY27 YIP Brief Topic Description\par}
{\small Prospective Program Officer: Dr. Behzad Kamgar-Parsi}
\end{center}
\section*{Overview}
Autonomous agents must learn unfamiliar facts, plan multi-step investigations, use tools, and revise conclusions as new observations arrive. Language models may complete a task with invalid intermediate steps or apply a correct rule under the wrong conditions. Training on final outcomes alone provides limited information about which reasoning behaviors are trustworthy and transferable.
\par\noindent\textbf{Central research question:} How can agents acquire unfamiliar external knowledge, learn checkable reasoning processes, and autonomously verify and revise conclusions?
\par We propose three connected research thrusts:
\begin{itemize}
\item \textbf{Structured knowledge for open-world reasoning.} We will retrieve and extract entities, events, relations, scientific principles, and contextual constraints from heterogeneous sources. Compact evidence records will distinguish observations, inferred relations, unresolved assumptions, and the conditions under which a scientific or logical rule applies.
\item \textbf{Scalable process-guided learning.} We will synthesize multi-hop reasoning tasks and intermediate trajectories from structured knowledge, including controlled reasoning errors. Fine-grained evaluators will identify factual, contextual, and logical failures. Process rewards will guide candidate-step selection and iterative training; we will test transfer to held-out knowledge compositions.
\item \textbf{Verification-driven autonomous reasoning.} Agents will form hypotheses or plans, expose required assumptions, select evidence-retrieval and symbolic or executable tools, and use diagnostic feedback to repair unsupported steps. Scientific hypothesis discovery, spanning literature retrieval and executable validation, will provide a challenging testbed for open-world generalization.
\end{itemize}
\noindent\textbf{Evaluation.} Evaluation will measure reasoning-step validity, outcome accuracy, error localization, repair success, transfer across conditions, and computational cost. Baselines will include retrieval-augmented agents, outcome-only learning, conventional process reward models, and tool-using agents without explicit condition tracking.
\par\noindent\textbf{Naval relevance and potential impact.} The project develops foundational methods for intelligent autonomous agents capable of reasoning under incomplete information and checking why a proposed action or conclusion is justified, with potential long-term relevance to Naval autonomy and planning.
\end{document}
